% v2.3.1 figure UID: FIG-P575-01
% Proposed post-v2.3.1 polish for FIG-P575-01: protect key-label size and stroke hierarchy.
\tikzset{slfig-FIG-P575-01/.style={
  every path/.append style={line cap=round,line join=round},
  every node/.style={font=\fontsize{9.2pt}{11pt}\selectfont}
}}
\begin{figure}[htbp]
\centering
\begin{tikzpicture}[slfig-FIG-P575-01]
\begin{groupplot}[group style={group size=2 by 1,horizontal sep=18mm},
  width=6.2cm,height=5.5cm,xmin=0,xmax=4,ymin=0,ymax=1,
  axis lines=left,xlabel={$x$},ylabel={$F(x)$},
  xtick={0,1,2,3,4},ytick={0,.25,.5,1},
  tick label style={font=\fontsize{8.5pt}{10.2pt}\selectfont},
  label style={font=\fontsize{9.2pt}{11pt}\selectfont},
  xlabel style={yshift=-8pt},
  title style={font=\fontsize{9.6pt}{11.5pt}\selectfont},
  axis line style={line width=.62pt},tick style={line width=.52pt},clip=false]
\nextgroupplot[title={连续：首次达到}]
  \addplot[draw=SLBlue,line width=1.0pt,domain=0:4,samples=201]{1-exp(-.65*x)};
  \draw[draw=SLTeal,line width=.65pt,densely dashed]
    (axis cs:0,.65)--(axis cs:1.615,.65)--(axis cs:1.615,0);
  \fill[SLTeal] (axis cs:1.615,.65) circle (2.2pt);
  \node[anchor=east,font=\fontsize{9.2pt}{11pt}\selectfont,text=SLTeal,xshift=-2pt]
    at (axis cs:0,.65) {$u=.65$};
  \node[anchor=north east,font=\fontsize{9.2pt}{11pt}\selectfont,text=SLTeal,
    xshift=-3pt,yshift=-8pt] at (axis cs:1.615,0) {$Q(.65)$};
\nextgroupplot[title={离散：首次达到}]
  \addplot[draw=SLBlue,line width=.95pt,const plot mark right]
    coordinates {(0,0)(1,.25)(2,.70)(3,.95)(4,1)};
  \foreach \x/\lo/\hi in {1/0/.25,2/.25/.70,3/.70/.95}{
    \addplot[only marks,mark=o,mark size=2.4pt,draw=SLBlue,fill=white] coordinates {(\x,\lo)};
    \addplot[only marks,mark=*,mark size=2.2pt,draw=SLBlue,fill=SLBlue] coordinates {(\x,\hi)};
  }
  \draw[draw=SLTeal,line width=.65pt,densely dashed]
    (axis cs:0,.70)--(axis cs:2,.70)--(axis cs:2,0);
  \draw[draw=SLGold,line width=.65pt,dash dot]
    (axis cs:0,.72)--(axis cs:3,.72)--(axis cs:3,0);
  \addplot[only marks,mark=square*,mark size=2.5pt,draw=SLTeal,fill=SLTeal] coordinates {(2,.70)};
  \addplot[only marks,mark=triangle*,mark size=2.7pt,draw=SLGold,fill=SLGold] coordinates {(3,.72)};
  \node[anchor=east,font=\fontsize{9.0pt}{10.8pt}\selectfont,text=SLTeal,
    xshift=-1pt,yshift=-6pt]
    at (axis cs:0,.70) {$u=.70$};
  \node[anchor=east,font=\fontsize{9.0pt}{10.8pt}\selectfont,text=SLGold,
    xshift=-1pt,yshift=8pt]
    at (axis cs:0,.72) {$u=.72$};
  \node[anchor=north east,font=\fontsize{9.0pt}{10.8pt}\selectfont,text=SLTeal,
    xshift=-4pt,yshift=-10pt] at (axis cs:2,0) {$Q(.70)=2$};
  \node[anchor=north west,font=\fontsize{9.0pt}{10.8pt}\selectfont,text=SLGold,
    xshift=4pt,yshift=-10pt] at (axis cs:3,0) {$Q(.72)=3$};
\end{groupplot}
\end{tikzpicture}
\caption{广义逆采用首次达到规则：连续情形沿水平线找到交点；离散情形中 $u=0.7$ 落在第二个跳点，而 $u=0.72$ 落在第三个跳点}\label{fig:V5-C02-generalized-inverse}
\end{figure}
